\section{Fundamentos y antecedentes}
\subsection{Contrato de cuantización y factorización}
Para una cuantización afín se interpreta un código entero mediante
\(\tilde r=s(q-z)\), con escala positiva \(s\) y zero-point entero \(z\)~\cite{jacob}.
Aquí se eligen pesos U4, \(q_w\in[0,15]\), y activaciones S8,
\(q_a\in[-128,127]\), con zero-point de activaciones igual a cero.
Para la fila \(r\), grupo \(g\) e índices \(I_g\), el subtotal es
\begin{equation}
D_{r,g}=\sum_{i\in I_g}(q_{w,r,i}-z_{w,r,g})q_{a,i}.
\label{eq:group}
\end{equation}
Con escalas compartidas dentro de cada grupo, una capa reconstruye
\begin{equation}
\tilde y_r=b_r+\sum_g s_{w,r,g}s_{a,g}D_{r,g}.
\end{equation}
El subtotal entero debe ser exacto; el error de cuantización y el del escalado
se evalúan por separado. Los grupos con escalas distintas no se combinan
como enteros para aplicar después una escala única arbitraria.

Por distributividad, la ecuación~\eqref{eq:group} también se calcula como
\begin{equation}
D_{r,g}=P_{r,g}-z_{w,r,g}S_{a,g},\qquad
P_{r,g}=\sum_{i\in I_g}q_{w,r,i}q_{a,i},\quad
S_{a,g}=\sum_{i\in I_g}q_{a,i}.
\label{eq:factored}
\end{equation}
Si varias filas consumen el mismo vector, \(S_{a,g}\) se reutiliza. El baseline
debe aprovechar esa propiedad. No hacerlo mediría también una desventaja
evitable de software.

\subsection{Antecedentes pertinentes y límite de la comparación}
XpulpNN estudia instrucciones SIMD de baja precisión y operaciones que
combinan producto punto y carga sobre procesadores RISC-V~\cite{xpulpnn}.
Dustin demuestra un contexto de aritmética de precisión mixta en un cluster
de varios cores~\cite{dustin}. Estos antecedentes impiden presentar el uso
de RISC-V o los productos punto de pocos bits como conceptos inéditos.

QServe conecta cuantización W4A8KV4 con kernels y costos de ejecución para
servicio de LLM en GPU~\cite{qserve}. Motiva considerar costos fuera de la
multiplicación, pero sus resultados no predicen el rendimiento de un core
RV32 docente. No se extrapolarán cifras de GPU, ASIC y FPGA entre plataformas.

\begin{tabularx}{\textwidth}{@{}lYY@{}}
\toprule
\textbf{Referencia} & \textbf{Relación con el proyecto} & \textbf{Lectura adicional requerida}\\\midrule
Jacob et al. & Cuantización y aritmética entera. & Compatibilidad del epílogo y reglas de redondeo.\\
XpulpNN & SIMD sub-byte y fusión de trabajo. & Semántica, operandos y soporte de correcciones.\\
Dustin & Precisión mixta en RISC-V. & Diferencias de ISA, microarquitectura y memoria.\\
QServe & W4A8 y costos de conversión. & Contratos de formatos, sin extrapolar speedups.\\\bottomrule
\end{tabularx}\par

La revisión es un punto de partida. Antes de enviar un artículo se ampliará
la comparación a operaciones empacadas y extensiones pertinentes, verificando
en las fuentes completas anchuras, signo, acumulación y transporte de metadatos.
La contribución definitiva se delimitará con esos antecedentes y los resultados.
